\documentclass[10pt,a4paper]{amsart}
\usepackage{amsmath,amssymb,amsthm,graphicx}
\usepackage{verbatim}
\newenvironment{reference}{\comment}{\endcomment}
%\newenvironment{reference}{}{}
\newenvironment{slogan}{\comment}{\endcomment}
\newenvironment{history}{\comment}{\endcomment}

% For commutative diagrams we use Xy-pic
\usepackage[all]{xy}

% We use 2cell for 2-commutative diagrams.
\xyoption{2cell}
% Theorem environments.
%
\theoremstyle{plain}
\newtheorem{theorem}[subsection]{Theorem}
\newtheorem{proposition}[subsection]{Proposition}
\newtheorem{lemma}[subsection]{Lemma}

\theoremstyle{definition}
\newtheorem{definition}[subsection]{Definition}
\newtheorem{example}[subsection]{Example}
\newtheorem{exercise}[subsection]{Exercise}
\newtheorem{situation}[subsection]{Situation}

\theoremstyle{remark}
\newtheorem{remark}[subsection]{Remark}
\newtheorem{remarks}[subsection]{Remarks}

\numberwithin{equation}{subsection}

% Macros
%
\def\lim{\mathop{\mathrm{lim}}\nolimits}
\def\colim{\mathop{\mathrm{colim}}\nolimits}
\def\Spec{\mathop{\mathrm{Spec}}}
\def\Hom{\mathop{\mathrm{Hom}}\nolimits}
\def\Ext{\mathop{\mathrm{Ext}}\nolimits}
\def\SheafHom{\mathop{\mathcal{H}\!\mathit{om}}\nolimits}
\def\SheafExt{\mathop{\mathcal{E}\!\mathit{xt}}\nolimits}
\def\Sch{\mathit{Sch}}
\def\Mor{\mathop{\mathrm{Mor}}\nolimits}
\def\Ob{\mathop{\mathrm{Ob}}\nolimits}
\def\Sh{\mathop{\mathit{Sh}}\nolimits}
\def\NL{\mathop{N\!L}\nolimits}
\def\CH{\mathop{\mathrm{CH}}\nolimits}
\def\proetale{{pro\text{-}\acute{e}tale}}
\def\etale{{\acute{e}tale}}
\def\QCoh{\mathit{QCoh}}
\def\Ker{\mathop{\mathrm{Ker}}}
\def\Im{\mathop{\mathrm{Im}}}
\def\Coker{\mathop{\mathrm{Coker}}}
\def\Coim{\mathop{\mathrm{Coim}}}

% Boxtimes
%
\DeclareMathSymbol{\boxtimes}{\mathbin}{AMSa}{"02}

%
% Macros for moduli stacks/spaces
%
\def\QCohstack{\mathcal{QC}\!\mathit{oh}}
\def\Cohstack{\mathcal{C}\!\mathit{oh}}
\def\Spacesstack{\mathcal{S}\!\mathit{paces}}
\def\Quotfunctor{\mathrm{Quot}}
\def\Hilbfunctor{\mathrm{Hilb}}
\def\Curvesstack{\mathcal{C}\!\mathit{urves}}
\def\Polarizedstack{\mathcal{P}\!\mathit{olarized}}
\def\Complexesstack{\mathcal{C}\!\mathit{omplexes}}
% \Pic is the operator that assigns to X its picard group, usage \Pic(X)
% \Picardstack_{X/B} denotes the Picard stack of X over B
% \Picardfunctor_{X/B} denotes the Picard functor of X over B
\def\Pic{\mathop{\mathrm{Pic}}\nolimits}
\def\Picardstack{\mathcal{P}\!\mathit{ic}}
\def\Picardfunctor{\mathrm{Pic}}
\def\Deformationcategory{\mathcal{D}\!\mathit{ef}}
\usepackage{cleveref}
\crefname{originallemma}{original lemma}{original lemmas}
\Crefname{originallemma}{Original lemma}{Original lemmas}
\makeatletter
\newlabel{sites-cohomology-lemma-cech-exact-presheaves}{{03AQ}{source}{}{sourcepointer-1}{}}
\newlabel{sites-cohomology-lemma-cech-exact-presheaves@cref}{{[originallemma][0][]03AQ}{[1][0][]source}{}{}{}}
\makeatother
\title{The Cech complex of a family of morphisms is an exact functor on abelian presheaves}
\author{The Stacks Project Authors}
\date{stacks-project:etale-cohomology.tex:local-snapshot-20261011; GNU FDL 1.2 or later}
\begin{document}
\maketitle
\noindent\textit{Editorial note.} All delivered mathematical text is the author's own, copied from the pinned Stacks Project chapter file etale-cohomology.tex; no statement or proof step was rewritten, reordered or removed.\par\medskip
\noindent\textit{Editorial note.} The proof is the author's complete argument: the pointwise criterion is stated, the degreewise product structure is identified, and the conclusion follows; the closing pointer to the same statement in another chapter is a reference and not a substitution for the argument.\par\medskip
\noindent\textit{Editorial note.} Every labelled reference occurring inside the counted spans is carried in assembly.external\_references as a non-load-bearing navigation pointer (same-chapter ones at their definition site, cross-chapter ones at the invocation inside the counted span), so no load-bearing mathematics is replaced by a stub. No bibliography citation occurs inside the counted spans.\par\medskip
\noindent\textit{Editorial note.} The checked chapter list for same-name collisions is the set of stacks chapters staged in this repository: chow, cohomology, divisors, etale-cohomology, flat, more-algebra, more-morphisms, morphisms, spaces-more-morphisms, varieties together with this card's chapter etale-cohomology.\par\medskip
\noindent\textit{Editorial note.} The preamble is the chapter's own native preamble excerpt already used for this source: the Xy-pic setup (native preamble lines 14-18) and the native theorem declarations with the mathematical macros (native preamble lines 154-217).\par\medskip
\noindent\textit{Allowed standard prerequisites (editorial).}\begin{itemize}
\item The Cech complex of a family of morphisms. Same chapter, Definition definition-cech-complex, recorded as the context of this card. Supplies the complex whose exactness is asserted.
\item Abelian presheaves form a functor category. Standard: the category of abelian presheaves on a category is the category of functors to abelian groups, hence abelian with exactness checked pointwise. Turns exactness of a sequence of presheaves into exactness of the sequences of values at every object.
\item The same statement in the chapter on cohomology on sites. Cohomology on Sites, Lemma sites-cohomology-lemma-cech-exact-presheaves. Is the pointer the author gives for the same computation.
\end{itemize}\medskip
\ifdefined\hypertarget\hypertarget{sourcepointer-1}{}\fi
\noindent\textit{Source reference.} Original lemma 03AQ: sources/stacks-etale-cohomology/pinned-original-source.tex, source lines 1530--1530. Navigation pointer invoked at proof line 1530; the counted span contains the pointer only and the cited result is not reproduced inside it, so the pointer is registered as a non-load-bearing navigation stub in the established form for Stacks material. No bibliography citation is involved.\par\medskip
\begin{definition}
\label{definition-cech-complex}
Let $\mathcal{C}$ be a category. Let
$\mathcal{U} = \{U_i \to U\}_{i \in I}$ be a family of morphisms of
$\mathcal{C}$ with fixed target. Assume that all the fibre products
$U_{i_0} \times_U \ldots \times_U U_{i_p}$ exists in $\mathcal{C}$.
Let $\mathcal{F} \in \textit{PAb}(\mathcal{C})$ be an abelian
presheaf. We define the {\it {\v C}ech complex}
$\check{\mathcal{C}}^\bullet(\mathcal{U}, \mathcal{F})$ by
$$
\prod_{i_0 \in I} \mathcal{F}(U_{i_0}) \to
\prod_{i_0, i_1 \in I} \mathcal{F}(U_{i_0} \times_U U_{i_1}) \to
\prod_{i_0, i_1, i_2 \in I}
\mathcal{F}(U_{i_0} \times_U U_{i_1} \times_U U_{i_2}) \to \ldots
$$
where the first term is in degree 0 and the maps are the usual ones.
The {\it {\v C}ech cohomology groups} are defined by
$$
\check{H}^p(\mathcal{U}, \mathcal{F}) =
H^p(\check{\mathcal{C}}^\bullet(\mathcal{U}, \mathcal{F})).
$$
\end{definition}

\begin{lemma}
\label{lemma-cech-presheaves}
Notation and assumptions as in Definition \ref{definition-cech-complex}.
The functor $\check{\mathcal{C}}^\bullet(\mathcal{U}, -)$
is exact on the category $\textit{PAb}(\mathcal{C})$.
\end{lemma}

\begin{proof}
This follows at once from the definition of a short exact sequence of
presheaves. Namely, as the category of abelian presheaves is the category of
functors on some category with values in $\textit{Ab}$, it is automatically an
abelian category: a sequence $\mathcal{F}_1\to \mathcal{F}_2\to \mathcal{F}_3$
is exact in $\textit{PAb}$ if and only if for all
$U \in \Ob(\mathcal{C})$, the sequence
$\mathcal{F}_1(U) \to \mathcal{F}_2(U) \to \mathcal{F}_3(U)$ is exact in
$\textit{Ab}$. So the complex above is merely a product of short exact
sequences in each degree. See also
Cohomology on Sites, Lemma \ref{sites-cohomology-lemma-cech-exact-presheaves}.
\end{proof}

\end{document}
